\documentclass[oneside]{uniaraxatcc} % Para impressão em frente e verso, utilize twoside
\usepackage{graphicx}
\usepackage{mathtext}
\usepackage{listings}
\usepackage{amssymb}
\usepackage{color}
\usepackage{subcaption}
\usepackage[utf8]{inputenc}
\usepackage{multirow}
\usepackage{scalefnt}
\usepackage{wrapfig}
\usepackage{nomencl}
\usepackage{lipsum}

\newcommand{\uniaraxa}{\uppercase{UNIARAXÁ}}

\tipotcc{P} % P:Projeto, M: Monografia
% Nome do Curso
\curso{Sistemas de Informação}

%Qual a habilitação do curso? (Bacharel, Tecnólogo ou Licenciado(a))
\habilitacao{Bacharel}

% Autor(a) do trabalho
\autor{Nome ALUNO1}

% Título do Trabalho
\titulo{Titulo do trabalho}

% Local cidade / estado
\local{Araxá}{Minas Gerais}


% Data dia / mês / ano da defesa
\date{04}{Outubro}{2022}

%Define a linha de pesquisa
% - Cultura, Desenvolvimento Humano e Gestão
% - Ambiente, Saúde e Políticas Públicas.
% - Produção, inovação e sustentabilidade.
\linhadepesquisa{Cultura e Desenvolvimento de Sistemas}

% Orientador [Tratamento]{Nome}{Gênero: M para Masculino e F para Feminino}
\orientador[Prof. Me]{Humberto Gustavo de Melo}{M}

% Membros da banca, podem ser definidos 2 membros de A a B (sendo que o orientador e o co-orientador já  são membros da banca)

% Primeiro membro da banca (Obrigatório)
\avaliadorA[Profa. Me.]{Nome do Membro 1}{F}

% Segundo membro da banca (Opcional)
\avaliadorB[Prof. Me.]{Nome do Membro 2 (Não obrigatório)}{M}

% Dedicatória (Opcional)
\dedicatoria{Dedicatória não obrigatória. }

% Agradecimentos (Opcional)
\agradecimentos{\noindent\chapter*{\uppercase{Agradecimentos}}

Texto de agradecimentos}

% Epígrafe (Opcional)
\epigrafe{\input{Partes/Epigrafe.tex}}

% Resumo (Obrigatório)
\resumo{\noindent\chapter*{\uppercase{Resumo}}

Resumo do trabalho do H1 aqui
  
\vspace{1cm}

\noindent \textbf{Palavras-chaves}:Jogos eletrônicos educativos, Análise de jogos educativos, Ensino-aprendizagem, Desenvolvimento de jogos educativos}

% Abstract (Obrigatório)
\abstract{\noindent\chapter*{\uppercase{Abstract}}

  \lipsum[1-3]

\vspace{1cm}

\noindent \textbf{Palavras-chaves}:Educational eletronic games, Educational games analysis, Teaching-learning, Development of educational games}

% Siglas e Abreviações (Opcional)
\siglas{\input{Partes/Siglas.tex}}


\begin{document}
\pretextual %Indica início dos elementos pré-textuais
\maketitle %Elementos Pré-textuais
\textual %Indica início dos elementos textuais
%\pagestyle{simple}

%\chapter{\uppercase{Tema}}
%\label{tema}
%Um estudo sobre como jogos digitais podem ser usados para ensinar crianças. %Utilizado somente no Projeto de TCC
\chapter{\uppercase{Introdução}}
\label{introducao}


 Segundo \citeonline{einstein}, é dito que nem todos tem a mesma velocidade de aprendizado, alguns aprendem com mais facilidade, e outros não. Essa dificuldade em aprender é complexa. Ela pode ter fatores comportamentais ou neurológicos, pode ter várias origens dentro ou fora da escola, podendo surgir em qualquer ponto da vida, e pode ser maior ainda se o aluno possuir transtornos como deficit de atenção, hiperatividade, dislexia ou falta de auto-estima. Deficiências físicas como problemas de visão e audição também podem atrapalhar, e algo simples como nascer com um cérebro que processa informações mais lentamente já pode impactar na aprendizagem de qualquer assunto e prática. O professor deve sempre levar em conta essa dificuldade, principalmente durante o ensino fundamental, onde as crianças estão aprendendo habilidades fundamentais para a vida adulta, como por exemplo, a matemática, uma disciplina fundamental, mas que porém, possui altos níveis de reprovação.

Assim como \citeonline{einstein} disse, as crianças aprendem e memorizam melhor brincando, ou se envolvendo em atividades estimulantes, como brincadeiras, onde podem colocar o seu conhecimento em prática enquanto se divertem, podendo errar e tentar novamente quantas vezes quiserem. As brincadeiras desempenham funções psicossociais, afetivas e intelectuais básicas no processo de desenvolvimento infantil, criando um ambiente propicio á aprendizagem que satisfaz a necessidade de brincar da criança e estimula a curiosidade e interações sociais de forma prazerosa.Uma outra forma existente de educar desta maneira são videogames, ou jogos eletrônicos. Com a popularidade de computadores e dispositivos eletrônicos de qualquer tipo no dia a dia, eles estão se tornando cada vez mais comuns na sala de aula, inclusive jogos eletrônicos, que são usados em sala de aula como uma ferramenta auxiliar de ensino já faz algum tempo, implementados de diversas formas para tornar a aprendizagem mais fácil e interessante pra os alunos, sendo especialmente úteis para os alunos que possuem dificuldade de aprender uma disciplina tão importante como matemática.

Neste Trabalho de Conclusão de Curso, o objetivo principal é mostrar, por meio de uma pesquisa descritiva, se usar jogos eletrônicos como uma ferramenta de ensino complementar no ensino de  matemática para crianças do ensino fundamental é justificável e como objetivo secundário criar um jogo 2D educativo sobre matemática que consiga ensinar sobre adição, subtração, a diferença entre números pares e ímpares e a existência de números negativos para crianças do ensino fundamental enquanto elas jogam um produto memorável e divertido.
\section{Problema e Justificativa}
\label{justificativa}

 \citeonline{einstein} diz que parte dos alunos possuem dificuldade em assimilar o conteúdo da disciplina de matemática, portanto, explicar os benefícios dos jogos digitais para a educação de crianças do ensino fundamental, e como os mesmos podem se tornar um modo de ensino opcional que atraia e prenda a atenção de seus potenciais alunos para que possa ensiná-los matemática enquanto se divertem é fundamental.


\section{Objetivos}
\label{objetivos}

%Assumindo que existe um problema a ser resolvido, apresente qual o objetivo de seu projeto de pesquisa. O que você pretende (ou pretendeu) exatamente fazer. Aqui, deve aparecer a principal ``contribuição'' de seu projeto. Qual é a principal ``coisa'' que você pretende/pretendeu fazer? Qual sua principal entrega? Não é necessário criar uma subseção para cada tipo. Pode haver uma única seção, chamada de ``objetivos'' cujo texto divida-se naturalmente em objetivo geral e objetivos específicos, deixando claro qual caso está sendo tratado em cada momento. Para diferenciar o objetivo geral dos objetivos específicos, siga as seguintes diretrizes:

\begin{itemize}
		\item \textbf{Objetivo geral}: Desenvolver um jogo educativo que auxilie o ensino de matemática matemática. 
		
		\item \textbf{Objetivos específicos}: Analisar se jogos eletrônicos devem ser usados em pró da educação infantil, mostrar exemplos de jogos educativos passados.
\end{itemize}

%\input{Partes/Hipotese.tex}
%\input{Partes/Limitacoes.tex}
%\input{Partes/Estrutura.tex}
\chapter{\uppercase{Referencial Teórico}}
\label{referencial}


\section{Seção 1}

 \lipsum[1-2]
 
 
\begin{figure}[h]
\centering
\captionsetup[subfigure]{labelformat=empty}
\caption{``Mascote H1''}
\begin{subfigure}{.5\textwidth}
\centering
\includegraphics[width=8cm,height=6cm]{Partes/Imagens/Mascote.png}
\caption{Fonte: Criada pelo autor(2019).}
\end{subfigure}%
\end{figure}
 
\section{Seção 2}
 
  \lipsum[1-9] 
%\input{Partes/Metodologia.tex}
%\input{Partes/Desenvolvimento.tex}
\chapter{\uppercase{Trabalhos Relacionados}}
\label{TrabalhosRelacionados}

 \lipsum[1-2] 
\chapter{\uppercase{Considerações finais}}
\label{conclusao}

Texto da conclusão aqui...


 \section{Trabalhos Futuros}
  
Citar o que pode ser feito no futuro deste trabalho

%\input{Partes/Etapas.tex} %Utilizado somente no Projeto de TCC
%\input{Partes/Cronograma.tex} %Utilizado somente no Projeto de TCC


\postextual %Indica início dos elementos pós-textuais	

    \bibliography{Referencias}

\begin{anexosenv}
% Imprime uma página indicando o início dos anexos
	%\partanexos

%(opcional)
\input{Partes/Anexo.tex} 
\end{anexosenv}


%(opcional)
\chapter{\uppercase{Considerações finais}}
\label{conclusao}

%Dicas para desenvolver a estrutura das considerações finais:

%\begin{itemize}

%\item Contextualize o projeto desenvolvido, ou seja, ajude o leitor a "relembrar" a estrutura. 

%\item Apresente os resultados.

%\item Apresente as contribuições do projeto para o meio acadêmico e/ou para o mercado. 

%\item Se for possível, descreva sobre melhorias para próximas pesquisas.

%\end{itemize}

%Neste Capítulo pode ser inserido as Seções: Trabalhos Futuros, Contribuições do Projeto, Artigos publicados. 

 Antes de iniciar este Trabalho de Conclusão de Curso, já havia diversos relatos falados vindos de conversas e palestras anteriores de como jogos eletrônicos eram importantes para outras atividades além de divertir, inclusive a pedagogia, na qual já se sabia que este conceito era experimentado faz alguns anos por conta da existência de jogos educativos. Ao desenvolver a pesquisa feita aqui, o objetivo inicial era confirmar estes relatos, usando a área de educação como exemplo, além de desenvolver o jogo educativo em paralelo mencionado anteriormente para treinar o desenvolvimento futuro de jogos e fazer um jogo educativo de matemática que pode possivelmente ser distribuído gratuitamente em escolas públicas.
 
 O desenvolvimento do jogo educativo ``Number Invasion" demorou, mas foi uma experiência enriquecedora. Foi um prazer trabalhar nele por conta da pretensão de seguir a área de desenvolvimento de jogos eletrônicos no futuro e adquirir experiência em um projeto como este era fundamental. Mais conhecimento sobre o Unity e a linguagem de programação C\# foi adquirido conforme o projeto foi se realizando, e diversas dúvidas sobre como alguns algoritmos eram feitos foram respondidas graças a enorme comunidade do Unity localizadas em fóruns e vídeos pela Internet, e ``Number Invasion" então se tornou o primeiro jogo completo feito, o que deu experiência valiosa e abre as portas para projetos futuros sem cunho educacional.
 
 A pesquisa foi a parte mais difícil deste Trabalho de Conclusão de Curso, pois explorava um campo novo onde não se havia experiência prévia como foi o desenvolvimento do jogo. Conforme a pesquisa se desenvolveu, o assunto foi aprendido aos poucos e o material da pesquisa foi complementando o jogo educativo e, com os dois unidos ao final da pesquisa, um completava o outro, onde a pesquisa mostrava a parte teórica e o jogo a parte prática.
 
 Com tudo o que foi apresentado ao decorrer deste Trabalho de Conclusão de Curso, conclui-se que os objetivos foram alcançados. O jogo educativo feito em paralelo, apesar de ter destino incerto quanto a ser distribuído em escolas, foi finalizado com sucesso sem apresentar erros e foi um bom treino para projetos futuros, e todas as evidências apresentadas de como o cérebro é afetado por jogos eletrônicos mostram que os mesmos têm grande potencial de serem usados como ferramentas auxiliares em diversas áreas, inclusive a pedagogia. Porém, é necessário estudar antes como será a implementação dos jogos eletrônicos em sala de aula, para que os mesmos sejam eficientes em sua função de ensinar e o máximo de conhecimento seja extraído.
 
  \section{Trabalhos Futuros}
  
  Este jogo não foi testado em sala de aula, portanto, sua capacidade de ensinar e ser divertido é desconhecida. Serão necessários testes práticos no futuro para comprovar sua eficácia. 

\end{document}
